%Key focus: Centrality meassures as stability/similarity measures
%Baseline that our models are good enough for our purpose as they double performance to a random model.
%CKA results and what that translates to in the graph stability
%Graph measures and expectations 
%Issues with bridging the gab between CKA and Graphs
%Issue witha benchmarking dataset
%Hidden chanal bharchart counts and chi squared
%Likelihood
%C3 and CP3 and the motter cortex
The focus of our project is to investigate potential similarities and stability within the learned graphs produced by our model. These graph representations aim to capture which nodes become salient during training, potentially revealing shared patterns that arise from learning the same underlying task across different initialisations.
\vspace{0.5em}\\
Before delving into the structure and consistency of these learned graphs, it is essential to first establish that our models are functionally valid. That is, that they perform well enough to justify further analysis. While an average test accuracy of 48.9\% may not be remarkable in absolute terms, it exceeds the 25\% baseline expected from random guessing in a 4 class classification task. This suggests that the models are learning non-trivial and task relevant features, providing a solid foundation for investigation.
\vspace{0.5em}\\
Importantly, our goal is not to maximize classification performance but to understand what the models are learning and whether those learned representations exhibit shared structure or stability across different training seeds. The doubled performance score relative to a random model is sufficient to support a discussion of internal graph representations, enabling us to explore how similar or different the resulting graphs are, and whether certain nodes or patterns consistently emerge as important across model instances.
\vspace{0.5em}\\
With this baseline established, the remainder of this discussion focuses on \sout{how we quantify representation} similarity using CKA, how we interpret the stability of graph metrics such as barycentre, and what these results reveal about the relationship between model learning dynamics and graph structures.
%CKA
\subsection{CKA Results}
When using CKA, we assess how similar the internal activations of two models are, layer by layer. This provides a high-level abstraction view of representational similarity without requiring exact alignment between individual neurons. CKA is useful for comparing models trained on the same task but with different random initialisations.
\vspace{0.5em}\\
In our analysis, the diagonal values of the CKA matrices are of primary interest. These compare the same layer across different models. High diagonal values indicate that different models tend to learn similar representations at corresponding depths, suggesting a degree of stability and shared internal structure. The off-diagonal values, which compare layers at different depths, are less relevant here, as our focus is on layer wise consistency rather than cross-layer relationships.
\vspace{0.5em}\\
On average, the top three layers yield CKA scores of over 0.9, while the final two remain relatively high, around 0.8. This consistency supports the assumption that the underlying learned graphs are likely to exhibit structural or functional similarities, though not exact equivalence. 
\vspace{0.5em}\\
\sout{The CKA requires data to be passed through the model in order to compute representation similarity. We use the test dataset for this purpose, as it offers a more reliable evaluation of model performance. One important note about our CKA implementation is that it includes a slight sampling bias: to ensure consistent and reasonable run times, we only evaluate the first half of the test }dataset (500 datapoints).
\subsection{CKA and Graph Metrics in Relation to Similarity}
While CKA provides a measure of representational similarity between models, it does not directly reflect the underlying graph structures produced by the DGCNN. Instead, CKA evaluates how similarly two models organise information in their internal activations, independent of specific neuron or node alignments. We performed CKA comparisons across the first 10 seeds, resulting in 90 pairwise similarity matrices. \sout{These were averaged into a single CKA matrix, which also showed high values along the diagonal, indicating that different model instances tend to learn similarly structured representations at corresponding layers. However, this similarity in activations does not necessarily imply structural similarity at the graph level, such as consistent importance of specific nodes or edges. To investigate that aspect, we turn to graph metrics such as graph centrality measures. }
\subsection{Graph measures and expectations}
To evaluate the stability of the extracted graph representations, we computed pairwise graph isomorphism among all graphs derived from the trained models. No two graphs were found to be isomorphic, indicating that each model produces a structurally distinct graph. This suggests a degree of instability in the \sout{learned representations}.
\vspace{0.5em}\\
Upon reflection, we recognise that while graph isomorphism is a sufficient condition to establish structural uniqueness, it is lacking. Isomorphism evaluates structural similarity while ignoring node labels. However, in our case, node labels correspond to EEG electrodes, and thus carry meaningful spatial and physiological significance. As a result, two graphs that are structurally identical but differ in node labelling would be functionally distinct.
\vspace{0.5em}\\
That said, this distinction does not affect our conclusions. Since none of the graphs are isomorphic, and isomorphism allows for node label permutation, it follows that none of the adjacency matrices are identical either.
\vspace{0.5em}\\
When calculating graph centrality metrics on the extracted graphs, we hypothesised that certain nodes would consistently emerge as central across models. This expectation was grounded in the fact that all models shared identical architectures and training dataset, with the main source of variability being the random initialisation of the trainable parameters.
\vspace{0.5em}\\
Our results largely support this hypothesis, with the electrodes C3, CP4, and CP3 emerging among the most frequent across all examined graph centrality measures. This includes barycentres, Laplacian centrality, betweenness centrality and current-flow betweenness centrality.
\vspace{0.5em}\\
Furthermore, comparisons between the trained and untrained model graphs show a clear shift in the central nodes. While in the untrained model graphs the distribution of barycentre counts is very similar to a uniform distribution, the trained model graphs show considerably higher counts for electrodes C3 and CP3 than for the other electrodes. These counts show significant increase from their expected values if the distribution was uniform. This difference is also highlighted by our Pearson's Chi squared test, which shows significant deviations from a uniform distribution at the 1\% significance level.

\subsection{Likelihood analysis}
The consistent emergence of C3 and CP3 as hubs across all centrality measures reveals properties of how DGCNNs process EEG data. These nodes exhibited a 9\% likelihood of being identified as centres, representing a 90\% increase over the uniform distribution baseline. This pattern stands in stark contrast to untrained models, which showed no meaningful structure with ~4.5\% centrality likelihood and less than 5\% deviation from uniform, demonstrating that the trained networks develop organised representations rather than random connectivity patterns. 

\subsection{Hidden channel bar chart counts and chi-squared analysis}
Given that we have found C3, CP3 and CP4 as the most common barycentres, we expected to see these electrodes as some of the most frequent nodes for other centrality measures as well.
We observed that for all of the centrality measures that we have experimented with, namely barycentres, betweenness centrality, Laplacian centrality and current-flow betweenness centrality, C3 and CP3 consistently appear as the two most frequent electrodes, followed by CP4.
\vspace{0.5em}\\
To further explore these patterns, we analysed centrality measure distributions grouped by the number of hidden channels in the corresponding model architectures. We examined bar charts grouped by each number of hidden channels listed in Table \ref{tab:hid_chans}, resulting in nine bar charts per centrality metric. Across all hidden channel configurations, C3, CP3, and CP4 remain consistently prominent, reinforcing the hypothesis that certain electrodes hold significance in the learned representations. Interestingly, the prominence of these electrodes appears to be clearer in models with a higher number of hidden channels, which could suggest that increased model capacity may amplify structural consistency. However, since the number of samples for each grouping is limited, further scaling of the experiment would be necessary to reach a conclusion. Additionally, we cannot validate these findings using another Pearson's Chi squared test due to some electrodes having a count lower than 5, which is the minimum count required for this statistical test. Increasing the number of samples or using an alternative statistical test could therefore be a consideration in future research.

\subsection{Issues with the benchmarking dataset}
In this project, we utilised the BNCI 2014-001 Motor Imagery dataset as the basis for training and evaluating the DGCNN models. The aforementioned dataset is widely adopted in the literature and serves as a common benchmark for EEG Motor Imagery tasks, making it a practical choice for our study. Despite this fact, the dataset presents certain limitations, such as it being a relatively old dataset, as well as the relative simplicity of the task.
\vspace{0.5em}\\
The experimental scope of this study was limited to a single dataset and a single model architecture, the DGCNN implementation provided by the \texttt{torcheeg} Python library. Despite the promising results of our study, an experiment at a bigger scale, comprising more datasets and more EEG tasks rather than just Motor Imagery would be required to asses whether learning consistent graph central nodes is a feature of this type of deep learning models. Only using this simple dataset has allowed us to focus on analysing graph stability, showcasing promising results which can be further explored and validated in future studies. 

\subsection{Physiological Alignment of Salient Nodes}
Another part of our analysis involves identifying nodes that consistently appear as central across multiple centrality measures in the learned graph structures. The electrodes C3, CP3, and CP4 consistently show the highest scores across all graphs in the Laplacian, betweenness, and current-flow betweenness centrality measures and in barycentre analyses. These measures were selected because they effectively capture different aspects of node importance in a graph. We hypothesise that the brain regions corresponding to these electrodes, being most active during motor imagery tasks, produce distinctive EEG patterns, which may lead to these nodes consistently appearing as central in the learned representations. This consistency suggests a stable and dominant role in the structure of the learned representations.
\vspace{0.5em}\\
According to the international 10-20 EEG electrode placement system and established neuroscientific literature, C3, CP3 and CP4 are situated over or near the motor cortex. This region of the brain is responsible for planning and executing voluntary movements. Specifically, C3 is located over the left motor cortex, which is the part of the brain responsible for controlling movement on the right side of the body. CP3 and CP4 are positioned just behind this region, over areas known as the sensorimotor cortex, whose function is coordinating movement and processing sensory information. 
\vspace{0.5em}\\
This is noteworthy in the context of our dataset, where participants were asked to imagine movements of their hands, feet, or tongue. The fact that these electrodes consistently emerge as central nodes in our learned graphs suggests that the DGCNN model instances are not only learning task-relevant features but are also capturing physiologically meaningful structures. This provides credibility to the model’s internal representations, as they appear to align with known brain anatomy. Moreover, the repeated emergence of the central-region nodes across different model initialisations further supports the stability of the learned representations. This consistency suggests that the networks are capturing reliable and interpretable patterns related to the motor imagery tasks. These findings point toward a connection between abstract model behaviour and concrete neurophysiological understanding, aiding both the interpretability and the trustworthiness of graph deep learning models in neuroimaging applications.
